\DomainHeader{Pipeline Development}{14\%}

\begin{Objectives}
  \item Convert USD to common formats and ensure fidelity.
  \item Document asset structure guidelines.
  \item Explain USDC vs USDA trade-offs.
  \item Implement round-trip pipelines between DCC and USD.
  \item Integrate custom resolvers to manage asset paths.
  \item Represent custom metadata.
  \item Validate asset paths are formatted correctly.
  \item Write or extend a USD importer in a DCC.
\end{Objectives}

\begin{Resources}
  \item Learn OpenUSD: \href{https://docs.nvidia.com/learn-openusd/latest/asset-structure/index.html}{Asset Structure Principles and Content Aggregation};
        \href{https://docs.nvidia.com/learn-openusd/latest/composition-basics/index.html}{Composition Basics};
        \href{https://docs.nvidia.com/learn-openusd/latest/beyond-basics/index.html}{Beyond the Basics}
  \item Technical references: \href{https://docs.omniverse.nvidia.com/usd/latest/learn-openusd/independent/asset-structure-principles.html}{Principles of Scalable Asset Structure};
        \href{https://github.com/usd-wg/assets/blob/main/docs/asset-structure-guidelines.md}{ASWF Asset Structure Guidelines};
        USD FAQ \href{https://openusd.org/release/usdfaq.html#i-have-some-layers-i-want-to-combine-should-i-use-sublayers-or-references}{``SubLayers or References?''}
  \item OpenUSD glossary: \href{https://openusd.org/release/glossary.html#usdglossary-subcomponent}{Subcomponent}; \href{https://openusd.org/release/glossary.html#crate-file-format}{Crate File Format}; \href{https://openusd.org/release/glossary.html#flatten}{Flatten}; \href{https://openusd.org/release/glossary.html#edittarget}{Edit Target}
  \item OpenUSD API: \href{https://openusd.org/release/api/class_usd_stage.html}{UsdStage::Flatten};
        \href{https://openusd.org/release/api/class_sdf_change_block.html}{SdfChangeBlock};
        \href{https://openusd.org/release/api/class_usd_notice.html}{UsdNotice}
\end{Resources}

\CheatSheetTitle
\begin{itemize}[leftmargin=*]
  \item \textbf{Asset interface pattern}: a small entry-point layer (stable
        name, \texttt{defaultPrim}, \texttt{kind}, \texttt{assetInfo})
        payloads the heavy contents. Consumers reference the interface and
        never see the internals.\index{asset structure}
  \item \textbf{Kinds nest}: \texttt{assembly} $>$ \texttt{group} $>$
        \texttt{component} $>$ \texttt{subcomponent}; models must form a
        \emph{contiguous} tree --- a component under a plain prim is not a
        model.\index{kind}\index{model hierarchy}
  \item \textbf{SubLayers vs References}: layers in the same stack share one
        namespace (departments on a shot); references graft subtrees under
        new paths (assets into scenes).
  \item \textbf{Custom metadata}: per-asset facts in \texttt{assetInfo};
        free-form notes in \texttt{customData}; neither is time-sampled.
  \item \textbf{Asset paths}: relative, forward slashes, resolver-friendly ---
        validate at publish time, not when a render farm machine fails.
  \item \textbf{Resolvers} (Ar) turn logical identifiers
        (\texttt{asset://chair}) into real files --- pipelines pin versions
        there, not in scene files.\index{resolver}
  \item \textbf{Delivery}: flatten (and \texttt{usdzip}) to strip composition
        and proprietary paths from outgoing assets.
\end{itemize}

\section{Core Concepts}

\subsection{The asset interface pattern}\label{sec:pd-iface}\objref{7.2}
Scalable pipelines standardize what an asset looks like \emph{from outside}:
one small entry-point layer with a stable prim name and metadata, which
payloads the actual contents. Publishing a new version means updating what
the interface points at --- no consumer scene changes.

\begin{lstlisting}
#usda 1.0
( defaultPrim = "Chair" )
def Xform "Chair" (
    kind = "component"
    assetInfo = {
        string name = "Chair"
        string version = "v003"
    }
    prepend payload = @./contents.usda@
)
{
}
\end{lstlisting}

\begin{lstlisting}
#usda 1.0
( defaultPrim = "Chair" )
def Xform "Chair" {
    def Scope "Geometry" {
        def Sphere "Seat" {
        }
    }
    def Scope "Looks" {
    }
}
\end{lstlisting}

\noindent\objref{7.6}Everything heavy lives behind the payload: consumers
who only need the scene skeleton load nothing, and
\texttt{assetInfo}\index{assetInfo} carries pipeline facts (name, version)
that tools can query without conventions-by-folklore.

\begin{figure}[ht]
  \centering
  \begin{tikzpicture}
    \node[layerbox, minimum width=2.6cm] (iface) {Chair.usda};
    \node[notetext, above=1.5mm of iface] {interface: defaultPrim, kind, assetInfo};
    \node[layerbox, minimum width=2.6cm, right=17mm of iface] (contents) {contents.usda};
    \node[primbox, minimum width=2.0cm, below left=8mm and -10mm of contents] (geo) {Geometry};
    \node[primbox, minimum width=2.0cm, below right=8mm and -10mm of contents] (looks) {Looks};
    \draw[payArc] (iface) -- node[arcLabel] {payload} (contents);
    \draw[arc, draw=black!50] (contents.south) -- (geo.north);
    \draw[arc, draw=black!50] (contents.south) -- (looks.north);
    \node[layerbox, minimum width=2.4cm, left=14mm of iface] (scene) {scene.usda};
    \draw[refArc] (scene) -- node[arcLabel] {reference} (iface);
  \end{tikzpicture}
  \caption{The asset interface pattern: scenes reference a thin, stable entry
  point; the entry point payloads the heavy, evolving contents.}
  \label{fig:asset-iface}
\end{figure}

\subsection{Kinds and the model hierarchy}\label{sec:pd-kinds}\objref{7.2, 3.3}
\index{kind}\index{model hierarchy}Kinds make large scenes navigable:
\texttt{assembly} (published groups of models), \texttt{group} (organizational
interior nodes), \texttt{component} (the leaf-level publishable asset),
\texttt{subcomponent} (named pieces inside a component). Tools rely on the
\emph{model hierarchy rule}: every ancestor of a model must itself be a group
(or assembly). Break the chain and queries silently stop seeing your asset:

\begin{lstlisting}
#usda 1.0
( defaultPrim = "City" )
def Xform "City" ( kind = "assembly" ) {
    def Xform "Block" ( kind = "group" ) {
        def Xform "Bldg" ( kind = "component" ) {
            def Xform "Door" ( kind = "subcomponent" ) {
            }
        }
    }
    def Xform "Plain" {
        def Xform "Orphan" ( kind = "component" ) {
        }
    }
}
\end{lstlisting}

\begin{lstlisting}[language=Python]
from pxr import Usd

stage = Usd.Stage.Open("hierarchy.usda")
for path in ("/City", "/City/Block", "/City/Block/Bldg",
             "/City/Block/Bldg/Door", "/City/Plain/Orphan"):
    prim = stage.GetPrimAtPath(path)
    print(path, Usd.ModelAPI(prim).GetKind(),
          "model" if prim.IsModel() else "-")
\end{lstlisting}

\begin{Gotcha}
Run it (the book did): \texttt{/City/Plain/Orphan} has
\texttt{kind = component} yet \texttt{IsModel()} is \textbf{False} --- its
parent \texttt{Plain} has no group kind, so the model chain is broken.
Model-hierarchy traversals and ``find all components'' tools will skip it.
The fix is organizational, not code: give interior prims group kinds.
\end{Gotcha}

\subsection{SubLayers or references?}\label{sec:pd-subref}\objref{1.3, 1.5}
The FAQ question is an exam question in disguise. \emph{Sublayers} stack whole
layers into one shared namespace --- right when several parties author
opinions about the \emph{same} prims (shot workflows: layout, animation,
lighting). \emph{References} graft a subtree under a chosen path --- right
when content must appear (possibly many times) \emph{inside} something else
(assets into sets, sets into shots). A useful tiebreaker: if you would need to
re-path prims to combine the files, you wanted a reference.

\subsection{Validating asset paths}\label{sec:pd-paths}\objref{7.7}
\index{asset path}Published layers travel: farm machines, other sites, other
studios. Absolute paths and backslashes are how assets die in transit. The
dependency list is one call away, so gate it at publish:

\begin{lstlisting}[language=Python]
from pxr import Sdf

layer = Sdf.Layer.FindOrOpen("Chair.usda")
for path in layer.GetCompositionAssetDependencies():
    ok = not path.startswith("/") and "\\" not in path
    print(("OK  " if ok else "BAD ") + path)
# prints: OK  ./contents.usda
\end{lstlisting}

\begin{ExamTip}
``Validate asset paths are formatted correctly'' $\Rightarrow$ relative to the
layer, forward slashes, resolvable by the studio resolver, and no machine-
or user-specific roots. Pair it with the export validator from the Data
Exchange chapter --- same gate, same publish step.
\end{ExamTip}

\subsection{Resolvers in the pipeline}\label{sec:pd-res}\objref{7.5, 3.5}
\index{resolver}An Ar resolver maps \emph{logical} asset identifiers to real
files at open time: \texttt{asset://props/chair} might resolve to
\texttt{/site/v003/chair.usdc} today and to a new version tomorrow ---
without touching a single scene file. That is where pinning, platform roots,
and site-specific storage belong. Scene files keep clean logical paths; the
resolver owns policy. (Writing one is plugin work --- see Customizing USD.)

\subsection{Delivery: flatten and strip}\label{sec:pd-deliver}\objref{1.9, 7.1}
To send an asset outside the pipeline, remove what outsiders cannot resolve:
\texttt{UsdStage::Flatten()} bakes composition (and resolves internal
references away); \texttt{usdzip} packages the result with its textures.
What you lose --- overridability, variants, payload deferral --- is exactly
what you do not want to leak.\index{flatten}\index{usdzip}

\begin{DeepDive}
Asset-structure \emph{guidelines} (objective: document them!) are a short
contract, not a wiki: entry-point naming, \texttt{defaultPrim} == asset name,
kind policy, where Geometry/Looks scopes live, primvar and namespace
conventions, path rules, and the validation gate that enforces all of it.
The ASWF Asset Structure Guidelines are the public reference point --- steal
their shape.
\end{DeepDive}

\DrillsTitle
\subsection*{Drill 1: Asset contract}
Write 10 rules for your ``publishable USD asset contract''.
\AnswerLines{10}
\begin{DrillSolution}
A solid contract: (1) entry-point layer named like the asset; (2)
\texttt{defaultPrim} set and equal to the root prim; (3) root prim
\texttt{kind = component} (or assembly for published groups); (4) heavy data
behind a payload; (5) \texttt{Geometry}/\texttt{Looks} scopes in fixed
places; (6) \texttt{assetInfo} carries name and version; (7) all asset paths
relative with forward slashes; (8) units and up-axis declared; (9) extents
authored on all Boundables; (10) passes the publish validator --- no
warnings, \texttt{usdchecker} clean.
\end{DrillSolution}

\subsection*{Drill 2: Naming and layout}
Propose a directory structure for an asset with geo, materials, and anim.
\AnswerLines{8}
\begin{DrillSolution}
One workable shape: \texttt{Chair/Chair.usda} (interface),
\texttt{Chair/contents.usda} (assembly of the parts), and
\texttt{Chair/parts/\{geometry,looks,anim\}.usda}. The interface payloads
\texttt{contents.usda}; contents sublayers (same namespace!) the parts.
Versioning happens by publishing new part files and re-pointing ---
consumers only ever reference \texttt{Chair/Chair.usda}.
\end{DrillSolution}

\subsection*{Drill 3: SubLayers or references?}
For each case pick one and justify: (a) lighting overrides on a shot; (b) a
chair placed 40 times in a restaurant; (c) merging two departments' edits to
the same character.
\AnswerLines{6}
\begin{DrillSolution}
(a) Sublayer --- same namespace, strength by list order. (b) References
(instanceable!) --- the same subtree grafted under 40 paths. (c) Sublayers
--- both edit the same prims; references would duplicate the character under
two roots.
\end{DrillSolution}

\section*{Checkpoint Questions}
\addcontentsline{toc}{section}{Checkpoint Questions}

\noindent\textbf{CQ1.} A tool that lists all components in a scene misses one
asset. Its kind is correctly \texttt{component}. What is the first thing to
check?
\begin{DrillSolution}
The ancestors: model hierarchy must be contiguous. If any ancestor lacks a
group/assembly kind, \texttt{IsModel()} is False for the component and
model-hierarchy traversals skip it --- verified in this chapter's listing.
\end{DrillSolution}

\noindent\textbf{CQ2.} Publishing v004 of an asset must not touch any of the
1{,}200 shots referencing it. Which two mechanisms make that possible?
\begin{DrillSolution}
The asset interface pattern (shots reference a stable entry-point layer, so
content behind it can change) and/or a resolver that maps a logical
identifier to the current published version. Both keep ``which version''
out of scene files.
\end{DrillSolution}

\noindent\textbf{CQ3.} Name the API call that lists a layer's external asset
dependencies for a path-validation gate.
\begin{DrillSolution}
\texttt{Sdf.Layer.GetCompositionAssetDependencies()} --- sublayer, reference,
and payload asset paths of that layer, ready to be checked for absolute
paths, backslashes, and resolvability.
\end{DrillSolution}

\subsection*{Mini-checklist (tick when mastered)}
\begin{itemize}[leftmargin=*]
  \item[$\square$] I can draw the asset interface pattern from memory.
  \item[$\square$] I can state the model-hierarchy contiguity rule and its failure mode.
  \item[$\square$] I can argue sublayers vs references for any concrete case.
  \item[$\square$] I can name the publish-time gates (paths, metadata, extents, checker).
\end{itemize}
